% Part: first-order-logic
% Chapter: proof-systems
% Section: axiomatic-deduction

\documentclass[../../../include/open-logic-section]{subfiles}

\begin{document}

\iftag{FOL}
      {\olfileid{fol}{prf}{axd}}
      {\olfileid{pl}{prf}{axd}}

\olsection{\usetoken{P}{derivation} به روش اصل‌موضوعی}

!!{derivation}sی اصل‌موضوعی، قدیمی‌ترین و ساده‌ترین دستگاه‌های منطقی برای
!!{derivation} هستند. !!{derivation}s آن‌ها صرفاً توالی‌هایی از
!!{sentence}s هستند. توالی‌ای از !!{sentence}s هنگامی !!^a{derivation}
درست به شمار می‌آید که هر !!{sentence}~$!A$ در آن یکی از شرایط زیر را
برآورده کند:
\begin{enumerate}
\item $!A$ یک اصل موضوع باشد، یا
\item $!A$ !!^a{element} از مجموعهٔ داده‌شدهٔ~$\Gamma$ از !!{sentence}s باشد، یا
\item $!A$ به‌واسطهٔ یک قاعدهٔ استنتاج توجیه شده باشد.
\end{enumerate}
برای آنکه $!A$ اصل موضوع باشد، باید صورت یکی از چند شمای ثابت
!!{sentence} را داشته باشد. مجموعه‌های بسیاری از شماهای اصل موضوع وجود دارند
که دستگاه !!{derivation} رضایت‌بخشی (صحیح و کامل) برای منطق مرتبهٔ اول فراهم
می‌کنند. برخی بر حسب رابط‌هایی که بر آن‌ها حاکم‌اند سازمان یافته‌اند؛ برای
نمونه، شِماهای
\[
!A \lif (!B \lif !A) \qquad !B \lif (!B \lor !C) \qquad (!B \land !C) \lif !B
\]
اصول موضوع متداولی‌اند که بر $\lif$، $\lor$ و~$\land$ حاکم‌اند. برخی
دستگاه‌های اصل‌موضوعی دستیابی به کمترین شمار اصول موضوع را هدف می‌گیرند.
بسته به رابط‌هایی که اولیه گرفته می‌شوند، حتی می‌توان دستگاه‌هایی
اصل‌موضوعی یافت که تنها از یک اصل موضوع تشکیل شده‌اند.

قاعدهٔ استنتاج گزاره‌ای شرطی است که شرطی کافی برای موجه‌بودن
!!^a{sentence} در !!^a{derivation} به دست می‌دهد. وضع مقدم یکی از قواعد
بسیار متداول از این نوع است: می‌گوید اگر $!A$ و $!A \lif
!B$ از پیش توجیه شده باشند، آنگاه $!B$ توجیه شده است. یعنی سطری از
!!^a{derivation} که !!{sentence}~$!B$ را دربر دارد توجیه شده است، مشروط
بر اینکه هم $!A$ و هم $!A \lif !B$ (برای یک !!{sentence}~$!A$) پیش
از~$!B$ در !!{derivation} ظاهر شده باشند.

رابطهٔ $\Proves$ مبتنی بر !!{derivation}sی اصل‌موضوعی چنین تعریف می‌شود:
$\Gamma \Proves !A$ اگر و تنها اگر !!^a{derivation}ی وجود داشته باشد که
!!{sentence}~$!A$ آخرین فرمول آن باشد (و $\Gamma$ مجموعهٔ !!{sentence}sیی
گرفته شود که در آن !!{derivation} با بند~(2) بالا توجیه شده‌اند). $!A$
قضیه است هرگاه~$!A$ !!^a{derivation}ی داشته باشد که در آن~$\Gamma$ تهی
است؛ یعنی هر !!{sentence} در !!{derivation} یا با بند (1) یا با بند~(3)
توجیه شده باشد. برای نمونه، !!^a{derivation} زیر نشان می‌دهد که $\Proves
!A  \lif (!B \lif (!B \lor !A))$:
\begin{derivation}
  1. & $!B \lif (!B \lor !A)$ \\
  2. & $(!B \lif (!B \lor !A)) \lif (!A  \lif (!B \lif (!B \lor !A)))$\\
  3. & $!A  \lif (!B \lif (!B \lor !A))$
\end{derivation}
!!{sentence} سطر~1 به صورت اصل موضوع $!A \lif (!A
\lor !B)$ است (با جابه‌جایی نقش‌های $!A$ و $!B$). جملهٔ سطر~2 به صورت
اصل موضوع $!A \lif (!B \lif !A)$ است. پس هر دو سطر توجیه شده‌اند. سطر~3
با وضع مقدم توجیه می‌شود: اگر آن را به صورت $!D$ کوتاه کنیم، آنگاه سطر~2
صورت $!C \lif !D$ را دارد، که در آن $!C$ همان
$!B \lif (!B \lor !A)$، یعنی سطر~1، است.

مجموعهٔ $\Gamma$ ناسازگار است هرگاه $\Gamma \Proves \lfalse$. هر دستگاه
اصل‌موضوعی کامل همچنین $\lfalse \lif !A$ را برای هر~$!A$ اثبات می‌کند؛
پس اگر $\Gamma$ ناسازگار باشد، آنگاه $\Gamma \Proves !A$ برای هر~$!A$.

دستگاه‌های منطقی با !!{derivation}sی اصل‌موضوعی را نخستین بار گوتلوب فرگه
در اثر سال ۱۸۷۹ خود، \emph{مفهوم‌نگاشت}، عرضه کرد؛ از همین رو، آن اثر را
اغلب نخستین اثر منطق جدید می‌دانند. این دستگاه‌ها در
\emph{اصول ریاضیات} آلفرد نورث وایتهد و برتراند راسل و نیز به دست
داوید هیلبرت و شاگردانش در دههٔ ۱۹۲۰ تکمیل شدند. ازاین‌رو، آن‌ها را اغلب
«دستگاه‌های فرگه» یا «دستگاه‌های هیلبرت» می‌نامند. این دستگاه‌ها از این
حیث بسیار انعطاف‌پذیرند که یافتن دستگاهی اصل‌موضوعی برای یک منطق غالباً
آسان است. چون !!{derivation}s ساختاری بسیار ساده و فقط یک یا دو قاعدهٔ
استنتاج دارند، اثبات مطالبی \emph{دربارهٔ} آن‌ها نیز نسبتاً آسان است.
با این حال، کاربرد عملی آن‌ها بسیار دشوار است؛ یعنی یافتن و نوشتن
برهان‌ها دشوار است.

\end{document}
